\chapter{Einschränken, Zusammenfügen und Erweitern}
\label{b05guide:ch:restrictions}

\section{Einschränkungen und Zielmengen}

\begin{definition}[Einschränkung einer Funktion]
\label{b05guide:def:restriction}
Für $f\colon A\to B$ und $C\subseteq A$ sei
\[
 f\restriction_C
   =\{(x,y)\in f\mid x\in C\}
   =f\cap(C\times B).
\]
Die Einschränkung behält die Werte an den Stellen aus $C$ bei.
\end{definition}

\begin{theorem}[Grundgesetze der Einschränkung]
\label{b05guide:thm:restriction-laws}
Seien $f\colon A\to B$ und $C,D\subseteq A$. Dann gilt
\[
 \begin{aligned}
 f\restriction_C&\colon C\to B,&
 (f\restriction_C)(x)&=f(x)\quad(x\in C),\\
 f\restriction_A&=f,& f\restriction_\varnothing&=\varnothing,\\
 (f\restriction_C)\restriction_{C\cap D}
    &=f\restriction_{C\cap D},\\
 f\restriction_{C\cup D}
    &=(f\restriction_C)\cup(f\restriction_D),\\
 f\restriction_{C\cap D}
    &=(f\restriction_C)\cap(f\restriction_D).
 \end{aligned}
\]
Für $U\subseteq C$ und $V\subseteq B$ gilt außerdem
\[
 (f\restriction_C)[U]=f[U],\qquad
 (f\restriction_C)^{-1}[V]=C\cap f^{-1}[V].
\]
\end{theorem}

\begin{theorem}[Genaues Kriterium für einen Wechsel der Zielmenge]
\label{b05guide:thm:restriction-target}
Seien $f\colon A\to B$, $C\subseteq A$ und $D$ eine beliebige Menge.
Für denselben eingeschränkten Graphen gilt
\[
 f\restriction_C\colon C\to D
 \quad\Longleftrightarrow\quad f[C]\subseteq D.
\]
Insbesondere ist $f\colon A\to D$ genau dann zulässig,
wenn $f[A]\subseteq D$. Für $D\subseteq B$ gilt ferner
\[
 f\cap(C\times D)
   =f\restriction_{C\cap f^{-1}[D]}
   \colon C\cap f^{-1}[D]\to D.
\]
Dieser Graph ist genau dann eine Funktion auf ganz $C$,
wenn $f[C]\subseteq D$.
\end{theorem}

\begin{theorem}[Teilgraphen und Fortsetzung]
\label{b05guide:thm:subgraphs}
Seien $f\colon A\to B$ und $H\subseteq f$. Mit
\[
 C=\{x\in A\mid\exists y\in B\,((x,y)\in H)\}
\]
ist $H=f\restriction_C$, also $H\colon C\to B$.
Das allgemeine Kriterium für enthaltene Funktionsgraphen steht in
Satz~\ref{b05guide:thm:graphinclusion}.
\end{theorem}

\section{Kompatible Graphen und Fallunterscheidungen}

\begin{theorem}[Vereinigung einer kompatiblen Familie]
\label{b05guide:thm:gluing}
Seien $I$ eine Menge und $f_i\colon A_i\to B$ für $i\in I$
vorgegebene Funktionen. Setze
\[
 A=\bigcup_{i\in I}A_i,\qquad H=\bigcup_{i\in I}f_i.
\]
Dann sind äquivalent:
\[
 \begin{aligned}
 &H\colon A\to B,\\
 &\forall i,j\in I\ \forall x\in A_i\cap A_j\,
                         (f_i(x)=f_j(x)).
 \end{aligned}
\]
Bei Erfüllung dieser Bedingung ist $H$ die eindeutig bestimmte
Funktion auf $A$ mit $H\restriction_{A_i}=f_i$ für alle $i\in I$.
Für $I=\varnothing$ sind $A=H=\varnothing$;
die Aussage gilt auch in diesem Fall.
\end{theorem}

\begin{theorem}[Durchschnitt von Funktionsgraphen]
\label{b05guide:thm:graph-intersection}
Sei $I\ne\varnothing$, und seien $f_i\colon A\to B$
für $i\in I$ vorgegeben. Die Menge der gemeinsamen
Übereinstimmungsstellen sei
\[
 D=\{x\in A\mid\forall i,j\in I\,(f_i(x)=f_j(x))\}.
\]
Für jedes $i_0\in I$ gilt
\[
 \bigcap_{i\in I}f_i=f_{i_0}\restriction_D\colon D\to B.
\]
Der Durchschnitt ist genau dann eine Funktion auf ganz $A$,
wenn alle $f_i$ gleich sind. Insbesondere ist für
$f,g\colon A\to B$ ihr Durchschnitt die Einschränkung
auf $\{x\in A\mid f(x)=g(x)\}$.
\end{theorem}

\begin{definition}[Fallabbildung und Überschreiben]
\label{b05guide:def:case-override}
Für $f\colon A\to B$ und $g\colon C\to B$ setze
\[
 \operatorname{Fall}(f,g)
       =f\cup(g\restriction_{C\setminus A}),\qquad
 \operatorname{Upd}(f,g)
       =(f\restriction_{A\setminus C})\cup g.
\]
Bei $\operatorname{Fall}(f,g)$ hat der erste Zweig Vorrang,
wie bei der Fallabbildung in Band~5.
Bei $\operatorname{Upd}(f,g)$ überschreibt $g$ die Werte
von $f$ auf dem gemeinsamen Definitionsbereich;
es gilt $\operatorname{Upd}(f,g)=\operatorname{Fall}(g,f)$.
\end{definition}

\begin{theorem}[Fallunterscheidung mit Vorrang]
\label{b05guide:thm:case-function}
Für $f\colon A\to B$ und $g\colon C\to B$ gilt
\[
 \begin{gathered}
 \operatorname{Fall}(f,g)\colon A\cup C\to B,\\
 \operatorname{Fall}(f,g)(x)=
 \begin{cases}
 f(x),&x\in A,\\
 g(x),&x\in C\setminus A.
 \end{cases}
 \end{gathered}
\]
Dabei gilt
\[
 \operatorname{Fall}(f,g)=f\cup g
 \quad\Longleftrightarrow\quad
 \forall x\in A\cap C\,(f(x)=g(x)).
\]
Ohne die letzte Bedingung ist der Fallgraph dennoch eine Funktion.
\end{theorem}

\begin{theorem}[Gesetze des Überschreibens]
\label{b05guide:thm:override}
Für $f\colon A\to B$ und $g\colon C\to B$ ist
$u=\operatorname{Upd}(f,g)$ eine Funktion von $A\cup C$ nach $B$.
Es gilt
\[
 u\restriction_C=g,\qquad
 u\restriction_{A\setminus C}=f\restriction_{A\setminus C}.
\]
Ist zusätzlich $h\colon D\to B$ gegeben, so gilt
\[
 \operatorname{Upd}(\operatorname{Upd}(f,g),h)
   =\operatorname{Upd}(f,\operatorname{Upd}(g,h)).
\]
Außerdem ist $\operatorname{Upd}(f,f)=f$;
das Überschreiben mit dem leeren Graphen verändert $f$ nicht.
\end{theorem}

\section{Erweiterungen und Änderung einzelner Werte}

\begin{theorem}[Erweiterung mit einem festen Wert]
\label{b05guide:thm:constant-extension}
Seien $f\colon A\to B$, $A\subseteq A'$ und $b\in B$.
Dann ist
\[
 h=f\cup\bigl((A'\setminus A)\times\{b\}\bigr)
       \colon A'\to B
\]
die eindeutig bestimmte Funktion mit
\[
 h\restriction_A=f,\qquad
 \forall x\in A'\setminus A\,(h(x)=b).
\]
Insbesondere ist $A'\times\{b\}$ eine konstante Funktion
auf $A'$. Für die Bildmenge der Erweiterung gilt
\[
 h[A']=
 \begin{cases}
 f[A],&A'=A,\\
 f[A]\cup\{b\},&A'\setminus A\ne\varnothing.
 \end{cases}
\]
\end{theorem}

\begin{theorem}[Genaues Existenzkriterium für eine Erweiterung]
\label{b05guide:thm:extension-existence}
Seien $f\colon A\to B$ und $A\subseteq A'$. Dann gilt
\[
 \begin{gathered}
 \exists h\colon A'\to B\ (h\restriction_A=f)\\
 \Longleftrightarrow\quad A'=A\ \lor\ B\ne\varnothing.
 \end{gathered}
\]
Die Aussage benötigt kein Auswahlaxiom.
\end{theorem}

\begin{theorem}[Hinzufügen und Ersetzen eines einzelnen Wertes]
\label{b05guide:thm:point-update}
Seien $f\colon A\to B$, $b\in B$ und $a$ beliebig.
Das bloße Hinzufügen des Paares $(a,b)$ erfüllt
\[
 f\cup\{(a,b)\}\colon A\cup\{a\}\to B
 \quad\Longleftrightarrow\quad
 a\notin A\ \lor\ (a,b)\in f.
\]
Unabhängig davon ist der ersetzende Graph
\[
 u=(f\restriction_{A\setminus\{a\}})\cup\{(a,b)\}
       \colon A\cup\{a\}\to B
\]
stets eine Funktion. Er ist durch
\[
 u(a)=b,\qquad
 \forall x\in A\setminus\{a\}\,(u(x)=f(x))
\]
eindeutig bestimmt. Für $a\notin A$ ist dies eine Erweiterung;
für $a\in A$ wird genau die vorgegebene Stelle überschrieben.
\end{theorem}
